\section{Introduction}
\label{sec:intro}

Congressional redistricting requires partitioning a state's census-tract
adjacency graph into $k$ contiguous, population-balanced districts.
Two algorithmic strategies dominate the graph-partitioning literature:
\emph{n-way} (or direct $k$-way) partitioning, which optimises all $k$ parts
simultaneously, and \emph{recursive bisection} (RB), which divides the graph
into two halves and recurses until $k$ leaves are reached.
Both strategies are implemented in METIS~\citep{karypis1998}, the
state-of-the-art multilevel graph partitioner, and both have been applied to
redistricting problems in the algorithmic literature~\citep{mehrotra1998optimal,fifield2020automated}.

The redistricting literature has characterised these two families primarily
for congressional chambers, where the district count $k$ is typically small
(1--52).  But redistricting practitioners face a far broader range of
configurations: state legislative chambers run from $k = 20$ (Nevada Senate)
to $k = 400$ (New Hampshire House), and the choice of partitioning strategy
may matter more at large $k$ where n-way partitioning's global optimisation
can be expected to shine.

Prior work in this portfolio (B.3, B.4) established that RB and n-way
yield statistically equivalent compactness for congressional-scale maps
under edge weighting.  However, those results were limited to a single
chamber type and a single geographic resolution (census tracts, $\sim$74{,}000
units nationally).  A practitioner implementing the Districting Integrity
Act (DIA) for multiple chamber types must know whether RB's compactness
advantage holds across state house and senate chambers, and whether the
runtime cost changes the answer at higher $k$.

\subsection*{Contributions}

This paper makes four contributions:

\begin{enumerate}
  \item We evaluate RB vs.\ n-way across three chamber types
    (congressional, state senate, state house) covering 435 congressional
    districts, 99 senate chambers, and 193 house chambers from the 2020
    Census, at both tract and block-group resolution.
  \item We show that RB's compactness advantage ($+0.003$--$+0.004$ mean
    Polsby--Popper) is consistent across all chamber types and persists at
    large $k$.
  \item We quantify the runtime crossover: n-way is 18--34\% faster only for
    state house chambers with $k > 80$, but the absolute difference is under
    200\,ms per state on a 16-core workstation.
  \item We identify the conditions under which RB's advantage narrows ---
    prime $k$ values that force asymmetric bisection trees --- and propose a
    practical mitigation.
\end{enumerate}

\subsection*{Paper structure}

Section~\ref{sec:related} reviews prior work on graph partitioning for
redistricting.  Section~\ref{sec:methodology} describes the experimental
design.  Section~\ref{sec:results} presents results by chamber type.
Section~\ref{sec:discussion} interprets the findings and connects them to
the DIA.  Section~\ref{sec:conclusion} concludes.

\subsection*{Why this comparison matters}

Algorithm selection is not politically neutral.  B.0 demonstrated a 12.8
percentage-point partisan gap between the unweighted baseline and GeoSection
for Wisconsin's 8 congressional districts.  If RB and n-way produced
systematically different partisan outcomes across chamber types, the DIA's
algorithm specification would face a legal challenge: choosing RB over n-way
would constitute an implicit partisan choice.  This paper shows that the
compactness difference between RB and n-way is small enough (and symmetric
in direction) that the choice is defensible on purely algorithmic grounds.
